% !TEX root = ./article/chapter1-article.tex

This chapter fixes notation and recalls standard definitions from
modal logic.

\section{The Modal Language}

We work in a modal language $\langm$ with variables
$\mathsf{Prop} = \set{p, q, r, \ldots}$, the Boolean connectives,
and a unary operator $\nec$ (necessity).  The dual
$\poss A \coloneqq \neg\nec\neg A$ is read ``possibly $A$.''

\begin{definition}\label{def:frame}
  A \textit{Kripke frame} is a pair $\FR = \oset{W, R}$
  where $W \neq \emptyset$ and $R \subseteq W \times W$.
\end{definition}

\begin{definition}\label{def:model}
  A \textit{Kripke model} $\MRel = \oset{W, R, V}$
  adds a valuation $V : \mathsf{Prop} \to \power{W}$ to a frame.
  The modal clause is:
  \[
    \MRel,w\trues \nec A \text{ if and only if }
    \text{for all } v,\ wRv \text{ implies } \MRel,v\trues A.
  \]
\end{definition}

\section{Normal Modal Logics}

The smallest normal modal logic, $\K$, extends classical
propositional logic with:
\begin{hilbertlist}
  \item[\kax] $\nec(A \to B) \to (\nec A \to \nec B)$
  \item[\nrule] From $\proves A$, infer $\proves \nec A$.
\end{hilbertlist}
Familiar extensions include $\KT = \K \oplus \tax$ (where $\tax$ is
$\nec A \to A$) and $\Sfour = \KT \oplus (\nec A \to \nec\nec A)$.

\begin{theorem}\label{thm:K-completeness}
  $\K$ is sound and complete with respect to the class of all
  Kripke frames \parencite[see][ch.~4]{blackburn2001}.
\end{theorem}

\begin{proof}
  Soundness: the axiom $\kax$ is valid in every frame and $\nrule$
  preserves validity.  Completeness: construct the canonical model
  $\Mlog = \oset{\Wlog, \Rlog, \Vlog}$ where
  $\Wlog$ is the set of maximal $\K$-consistent sets and
  $\Gamma \Rlog \Delta$ iff $\set{A : \nec A \in \Gamma}
  \subseteq \Delta$.  A routine induction (the Truth Lemma) shows
  every consistent formula is satisfiable in $\Mlog$.
\end{proof}
